\section{Banco de questões — Auditoria do Setor Público}

\subsection*{N1 — Fundamento competitivo}

\begin{questao}{\nivel{N1} 19.01 — Cebraspe/Certo ou Errado}
Critério de auditoria é a referência usada para avaliar o objeto; condição é a situação efetivamente encontrada pelo auditor.
\end{questao}

\begin{questao}{\nivel{N1} 19.02 — Cebraspe/Certo ou Errado}
Ceticismo profissional significa presumir fraude por parte do auditado até que ele prove o contrário.
\end{questao}

\begin{questao}{\nivel{N1} 19.03 — Cebraspe/Certo ou Errado}
Julgamento profissional deve ser fundamentado e documentado; não é autorização para substituir critérios por preferência pessoal do auditor.
\end{questao}

\begin{questao}{\nivel{N1} 19.04 — Cebraspe/Certo ou Errado}
Materialidade pode possuir dimensão qualitativa; uma falha de baixo valor financeiro pode ser relevante por envolver dados sigilosos, fraude ou serviço essencial.
\end{questao}

\begin{questao}{\nivel{N1} 19.05 — Cebraspe/Certo ou Errado}
Suficiência refere-se predominantemente à quantidade de evidência, enquanto apropriação envolve sua relevância e confiabilidade.
\end{questao}

\begin{questao}{\nivel{N1} 19.06 — Cebraspe/Certo ou Errado}
Indagação isolada ao gestor é, em regra, evidência suficiente para concluir sobre a efetividade de controle crítico de segurança.
\end{questao}

\begin{questao}{\nivel{N1} 19.07 — Cebraspe/Certo ou Errado}
Teste de controle busca avaliar desenho e/ou operação de controles; procedimento substantivo busca detectar diretamente distorções ou irregularidades nos dados, transações ou saldos.
\end{questao}

\begin{questao}{\nivel{N1} 19.08 — Cebraspe/Certo ou Errado}
Monitoramento de deliberações se encerra quando o gestor informa que iniciou providências, ainda que não exista evidência de implementação ou efeito.
\end{questao}

\begin{questao}{\nivel{N1} 19.09 — Cebraspe/Certo ou Errado}
Um achado de auditoria robusto tende a relacionar critério, condição, causa, efeito/risco e evidência, além de encaminhamento compatível.
\end{questao}

\begin{questao}{\nivel{N1} 19.10 — Cebraspe/Certo ou Errado}
Uma recomendação excessivamente prescritiva pode transferir ao auditor decisão gerencial que deveria permanecer com a administração.
\end{questao}

\subsection*{N2 — Aplicação}

\begin{questao}{\nivel{N2} 19.11 — FGV/tipos de auditoria}
Uma equipe avalia se programa de transformação digital alcançou resultados previstos, com uso econômico de recursos e boa relação entre insumos e produtos. O trabalho se aproxima de auditoria:
\begin{enumerate}[label=\Alph*)]
\item financeira;
\item operacional;
\item exclusivamente de conformidade;
\item de certificação digital;
\item contábil privada apenas.
\end{enumerate}
\end{questao}

\begin{questao}{\nivel{N2} 19.12 — FCC/3 Es}
Um data center entrega o mesmo volume de serviços com 25\% menos recursos, mantendo qualidade. A melhoria está mais diretamente ligada a:
\begin{enumerate}[label=\Alph*)]
\item efetividade;
\item eficiência;
\item legalidade apenas;
\item materialidade;
\item independência.
\end{enumerate}
\end{questao}

\begin{questao}{\nivel{N2} 19.13 — Cebraspe/Certo ou Errado}
Comprar pelo menor preço não demonstra, por si só, economicidade se o recurso adquirido tiver qualidade inadequada e elevar custos posteriores.
\end{questao}

\begin{questao}{\nivel{N2} 19.14 — FGV/critérios}
A questão de auditoria é: ``os acessos privilegiados são revisados trimestralmente?''. O critério mais apropriado é:
\begin{enumerate}[label=\Alph*)]
\item opinião do auditor sobre segurança;
\item política/norma aplicável que estabeleça periodicidade, papéis e requisitos da revisão;
\item quantidade de usuários do sistema;
\item custo do software de IAM;
\item relatório do fornecedor sem referência normativa.
\end{enumerate}
\end{questao}

\begin{questao}{\nivel{N2} 19.15 — Cebraspe/Certo ou Errado}
Definir objetivo, escopo e questões de auditoria antes de selecionar procedimentos ajuda a evitar coleta de evidência desconectada da conclusão pretendida.
\end{questao}

\begin{questao}{\nivel{N2} 19.16 — FCC/evidência}
O gestor afirma que backups são testados mensalmente. Para concluir sobre efetividade, o procedimento mais adequado é:
\begin{enumerate}[label=\Alph*)]
\item aceitar a declaração;
\item obter registros dos testes, selecionar amostras e verificar resultados/restaurações;
\item entrevistar novamente o mesmo gestor;
\item verificar apenas se existe software de backup;
\item considerar a existência de RAID como confirmação.
\end{enumerate}
\end{questao}

\begin{questao}{\nivel{N2} 19.17 — FGV/procedimentos}
O auditor recalcula o índice de disponibilidade a partir dos logs brutos e da fórmula contratual. O procedimento é principalmente:
\begin{enumerate}[label=\Alph*)]
\item indagação;
\item recálculo;
\item observação;
\item confirmação externa apenas;
\item inspeção física.
\end{enumerate}
\end{questao}

\begin{questao}{\nivel{N2} 19.18 — Cebraspe/Certo ou Errado}
Reexecução difere de recálculo porque pode envolver o auditor executar novamente um controle ou procedimento originalmente realizado pela entidade.
\end{questao}

\begin{questao}{\nivel{N2} 19.19 — FCC/amostragem}
Em amostragem estatística, o risco de amostragem corresponde à possibilidade de:
\begin{enumerate}[label=\Alph*)]
\item o auditor selecionar exatamente toda a população;
\item a conclusão baseada na amostra diferir da que seria obtida examinando toda a população;
\item qualquer evidência ser falsa;
\item o critério ser inadequado;
\item o relatório ser publicado fora do prazo.
\end{enumerate}
\end{questao}

\begin{questao}{\nivel{N2} 19.20 — FGV/materialidade}
Um sistema de denúncias custa R\$ 18 mil, mas expõe publicamente identidade de denunciantes. O auditor decide tratar o achado como material porque:
\begin{enumerate}[label=\Alph*)]
\item materialidade é sempre proporcional ao custo do ativo;
\item aspectos qualitativos de sigilo, direitos e impacto institucional podem ser relevantes;
\item sistemas baratos não podem ser materiais;
\item somente perda financeira importa;
\item materialidade não se aplica a TI.
\end{enumerate}
\end{questao}

\begin{questao}{\nivel{N2} 19.21 — Cebraspe/Certo ou Errado}
Risco de detecção pode ser reduzido pelo auditor ao alterar natureza, época e extensão dos procedimentos, enquanto risco inerente e risco de controle decorrem mais diretamente do objeto e dos controles da entidade.
\end{questao}

\begin{questao}{\nivel{N2} 19.22 — FCC/matriz de planejamento}
Qual combinação pertence naturalmente a uma matriz de planejamento?
\begin{enumerate}[label=\Alph*)]
\item questão, critério, informação necessária e procedimento;
\item somente achado e recomendação;
\item apenas cronograma e equipe;
\item nota fiscal e sanção;
\item conclusão sem evidência.
\end{enumerate}
\end{questao}

\begin{questao}{\nivel{N2} 19.23 — FGV/achado}
O auditor encontrou 7 de 18 contas administrativas compartilhadas, contrariando política que exige contas nominativas. A frase ``o controle de acesso é ruim'' é inadequada porque:
\begin{enumerate}[label=\Alph*)]
\item é ampla e não explicita critério, condição e evidência;
\item todo achado deve conter valor financeiro;
\item contas compartilhadas são permitidas sempre;
\item achados não podem tratar segurança;
\item a política nunca pode ser critério.
\end{enumerate}
\end{questao}

\begin{questao}{\nivel{N2} 19.24 — Cebraspe/Certo ou Errado}
Papéis de trabalho devem permitir que auditor experiente, sem participação prévia, compreenda objetivo, procedimentos, evidências, resultados e conclusões relevantes.
\end{questao}

\begin{questao}{\nivel{N2} 19.25 — FCC/monitoramento}
Uma recomendação foi classificada como implementada porque o órgão publicou nova política, mas nenhuma evidência mostra que as contas antigas foram revisadas. A avaliação mais adequada é:
\begin{enumerate}[label=\Alph*)]
\item a publicação da política prova a mitigação integral;
\item é necessário verificar implementação e, quando relevante, efetividade do controle;
\item política e execução são sempre equivalentes;
\item monitoramento deve se limitar a documentos normativos;
\item recomendações nunca exigem evidência posterior.
\end{enumerate}
\end{questao}

\subsection*{N3 — Integração de concurso}

\begin{questao}{\nivel{N3} 19.26 — FGV/assertivas ISSAI/NBASP}
Considere:

I. A auditoria do setor público envolve auditor, parte responsável e usuários previstos.

II. Critérios devem ser adequados ao objeto e à finalidade do trabalho.

III. Evidência suficiente e apropriada sustenta a confiança nas conclusões.

Assinale a alternativa correta.
\begin{enumerate}[label=\Alph*)]
\item Apenas I.
\item Apenas I e II.
\item Apenas II e III.
\item I, II e III.
\item Apenas III.
\end{enumerate}
\end{questao}

\begin{questao}{\nivel{N3} 19.27 — Cebraspe/Certo ou Errado}
Em trabalho de relatório direto, o auditor mede ou avalia o objeto segundo critérios, diferentemente da certificação em que a parte responsável apresenta a informação do objeto a ser auditada.
\end{questao}

\begin{questao}{\nivel{N3} 19.28 — FCC/auditoria operacional}
Um programa entrega 10 mil serviços digitais com custo unitário menor, mas não melhora o resultado público que motivou sua criação. O cenário indica:
\begin{enumerate}[label=\Alph*)]
\item possível ganho de eficiência sem demonstração de efetividade;
\item efetividade plena;
\item economicidade e efetividade são sinônimos;
\item ausência de output;
\item auditoria financeira apenas.
\end{enumerate}
\end{questao}

\begin{questao}{\nivel{N3} 19.29 — FGV/risco de auditoria}
O auditor avalia risco inerente e de controle como elevados. Mantido o risco de auditoria desejado em nível baixo, a resposta mais coerente é:
\begin{enumerate}[label=\Alph*)]
\item aceitar risco de detecção maior;
\item reduzir risco de detecção planejado por procedimentos mais robustos;
\item eliminar testes;
\item reduzir materialidade obrigatoriamente a zero;
\item confiar apenas em indagações.
\end{enumerate}
\end{questao}

\begin{questao}{\nivel{N3} 19.30 — Cebraspe/Certo ou Errado}
O modelo $RA=RI\times RC\times RD$ é conceitual e não exige que o auditor atribua probabilidades matematicamente exatas a cada componente para que seja útil no planejamento.
\end{questao}

\begin{questao}{\nivel{N3} 19.31 — FCC/evidência digital}
Um screenshot mostra que uma configuração estava correta no momento da captura. Qual limitação permanece?
\begin{enumerate}[label=\Alph*)]
\item screenshot prova configuração durante todo o período;
\item pode não demonstrar completude, período, origem ou impossibilidade de manipulação;
\item screenshot nunca é evidência;
\item imagem digital não pode ser documentada;
\item a evidência é automaticamente externa.
\end{enumerate}
\end{questao}

\begin{questao}{\nivel{N3} 19.32 — FGV/logs}
Logs usados para provar segregação podem ser apagados pelos mesmos administradores auditados. A implicação é:
\begin{enumerate}[label=\Alph*)]
\item confiabilidade da evidência é fragilizada e exige corroboração/controle de integridade;
\item os logs tornam-se prova absoluta;
\item basta aumentar a retenção local;
\item não há impacto se houver senha;
\item auditoria deve ignorar logs.
\end{enumerate}
\end{questao}

\begin{questao}{\nivel{N3} 19.33 — Cebraspe/Certo ou Errado}
Hash calculado sobre cópia de evidência digital pode demonstrar integridade da cópia ao longo da cadeia, mas não prova que o conteúdo original era verdadeiro.
\end{questao}

\begin{questao}{\nivel{N3} 19.34 — FCC/amostragem}
Em teste de controles, a taxa esperada de desvio aumenta significativamente. Mantidos os demais fatores, o tamanho de amostra tende a:
\begin{enumerate}[label=\Alph*)]
\item diminuir sempre;
\item aumentar para sustentar o nível de confiança desejado;
\item tornar-se zero;
\item ser independente da expectativa de desvio;
\item igualar-se obrigatoriamente à população.
\end{enumerate}
\end{questao}

\begin{questao}{\nivel{N3} 19.35 — FGV/causa do achado}
Auditoria encontra acessos de ex-servidores ativos. A chefia afirma que ``o problema é erro humano''. Antes de registrar essa causa, o auditor deve:
\begin{enumerate}[label=\Alph*)]
\item aceitá-la por ser plausível;
\item investigar processo de desligamento, integração RH-IAM, responsabilidades e controles;
\item excluir causa do achado;
\item atribuir culpa ao usuário;
\item recomendar apenas troca de senha.
\end{enumerate}
\end{questao}

\begin{questao}{\nivel{N3} 19.36 — Cebraspe/Certo ou Errado}
Causa mal identificada pode produzir recomendação que trata apenas o sintoma e não reduz a recorrência do achado.
\end{questao}

\begin{questao}{\nivel{N3} 19.37 — FCC/contraditório}
A unidade auditada apresenta evidência nova que invalida parte da condição descrita. O auditor deve:
\begin{enumerate}[label=\Alph*)]
\item ignorar porque o relatório já foi redigido;
\item avaliar a evidência e ajustar conclusão/achado se necessário;
\item apenas anexar sem análise;
\item manter o texto para preservar independência;
\item encerrar a auditoria automaticamente.
\end{enumerate}
\end{questao}

\begin{questao}{\nivel{N3} 19.38 — FGV/recomendação}
O achado identifica ausência de processo de recertificação de acessos. Qual recomendação é melhor?
\begin{enumerate}[label=\Alph*)]
\item ``compre a ferramenta X'';
\item ``implante processo periódico de recertificação, com responsáveis, critérios, evidência e tratamento de exceções'';
\item ``troque toda a equipe'';
\item ``garanta que nunca mais ocorra falha'';
\item ``aumente a segurança''.
\end{enumerate}
\end{questao}

\begin{questao}{\nivel{N3} 19.39 — Cebraspe/Certo ou Errado}
Recomendação orientada a resultado e causa tende a preservar melhor a responsabilidade gerencial do que prescrever produto específico sem demonstrar que seja a única solução viável.
\end{questao}

\begin{questao}{\nivel{N3} 19.40 — FCC/instrumentos de fiscalização}
O objetivo principal é conhecer arquitetura, processos e riscos de uma área ainda pouco compreendida, para subsidiar fiscalizações futuras. O instrumento mais aderente é:
\begin{enumerate}[label=\Alph*)]
\item levantamento;
\item monitoramento;
\item sanção;
\item recebimento definitivo;
\item auditoria financeira obrigatoriamente.
\end{enumerate}
\end{questao}

\subsection*{N4 — Auditoria / avançado}

\begin{questao}{\nivel{N4} 19.41 — Cebraspe/caso integrado}
Auditoria conclui que ``a empresa presta serviço ruim'' com base em três reclamações, sem critério contratual, universo, amostragem, SLA ou análise de chamados. Julgue o item:

A conclusão é ampla em relação à evidência disponível e carece de critério e procedimentos suficientes para sustentar a generalização.
\end{questao}

\begin{questao}{\nivel{N4} 19.42 — FGV/evidência de controle}
A política exige revisão trimestral de privilégios. O órgão apresenta planilha assinada pelo gestor com a frase ``revisado''. O auditor deseja concluir sobre operação efetiva. O melhor procedimento é:
\begin{enumerate}[label=\Alph*)]
\item aceitar a assinatura;
\item selecionar ciclos/amostras e verificar população, decisões, revogações e evidências do processo;
\item verificar apenas data do arquivo;
\item entrevistar o gestor e encerrar;
\item verificar se a ferramenta de IAM está licenciada.
\end{enumerate}
\end{questao}

\begin{questao}{\nivel{N4} 19.43 — FCC/materialidade}
Uma falha expõe dados sigilosos de denunciantes, mas não gera perda financeira mensurável. A melhor avaliação é:
\begin{enumerate}[label=\Alph*)]
\item imaterial por ausência de valor monetário;
\item potencialmente material por aspectos qualitativos e impacto institucional;
\item material apenas se atingir 5\% do orçamento;
\item irrelevante em auditoria operacional;
\item materialidade só se aplica a demonstrações financeiras.
\end{enumerate}
\end{questao}

\begin{questao}{\nivel{N4} 19.44 — Cebraspe/Certo ou Errado}
Uma análise de dados pode selecionar transações atípicas para exame, mas o score analítico não substitui evidência suficiente e apropriada para concluir pela irregularidade de cada caso.
\end{questao}

\begin{questao}{\nivel{N4} 19.45 — FGV/papéis de trabalho}
Um script SQL utilizado na auditoria é alterado após a execução, e a versão original não foi preservada. O risco principal é:
\begin{enumerate}[label=\Alph*)]
\item perda de reprodutibilidade do procedimento e da evidência obtida;
\item excesso de documentação;
\item impossibilidade de usar SQL;
\item falta de materialidade;
\item perda automática de independência institucional.
\end{enumerate}
\end{questao}

\begin{questao}{\nivel{N4} 19.46 — Cebraspe/Certo ou Errado}
Em auditoria de TI, preservar versão do código analítico, parâmetros e referência dos dados é parte importante da documentação quando esses elementos sustentam a conclusão.
\end{questao}

\begin{questao}{\nivel{N4} 19.47 — FCC/monitoramento}
Uma recomendação determinava revogar contas de desligados e implantar controle preventivo. O monitoramento comprova que contas antigas foram removidas, mas novos desligamentos continuam gerando contas órfãs. A conclusão é:
\begin{enumerate}[label=\Alph*)]
\item implementação e efetividade foram plenamente demonstradas;
\item houve correção pontual, mas o controle preventivo não se mostrou efetivo;
\item a recomendação deve ser considerada encerrada apenas porque o estoque foi zerado;
\item monitoramento não avalia recorrência;
\item o achado original era inválido.
\end{enumerate}
\end{questao}

\begin{questao}{\nivel{N4} 19.48 — FGV/independência}
O auditor recomenda produto específico e posteriormente participa da gestão da implantação do mesmo produto. O risco principal é:
\begin{enumerate}[label=\Alph*)]
\item assumir responsabilidade gerencial e comprometer objetividade em avaliação futura;
\item aumentar independência;
\item reduzir materialidade;
\item melhorar segregação;
\item eliminar risco de auditoria.
\end{enumerate}
\end{questao}

\begin{questao}{\nivel{N4} 19.49 — Cebraspe/Certo ou Errado}
A recomendação pode deixar à gestão a escolha da solução específica quando múltiplas alternativas são capazes de tratar a causa e atingir o resultado de controle esperado.
\end{questao}

\begin{questao}{\nivel{N4} 19.50 — FGV/matriz de achado}
Foram identificadas 42 contas de ex-servidores ativas por até 11 meses. Qual formulação de achado é mais adequada?
\begin{enumerate}[label=\Alph*)]
\item ``O Active Directory é inseguro'';
\item ``Os ex-servidores são responsáveis pela falha'';
\item ``A ausência de integração eficaz entre desligamento e revogação de acessos resultou em contas sem necessidade, ampliando risco de uso indevido; devem ser investigadas causas, privilégios e controles de offboarding'';
\item ``Basta trocar todas as senhas'';
\item ``Como não houve incidente conhecido, não existe efeito ou risco''.
\end{enumerate}
\end{questao}
